\documentclass[11pt]{article}

\usepackage[margin=1in]{geometry}
\usepackage{amsmath,amssymb,amsthm}
\usepackage{listings}
\usepackage{xcolor}
\usepackage{hyperref}
\usepackage{enumitem}
\usepackage{mathtools}

\newtheorem{proposition}{Proposition}
\newtheorem{lemma}{Lemma}
\newtheorem{corollary}{Corollary}
\newtheorem{example}{Example}

\definecolor{codebg}{rgb}{0.96,0.96,0.96}
\definecolor{codekw}{rgb}{0.0,0.0,0.6}
\definecolor{codecomment}{rgb}{0.35,0.45,0.35}

\lstdefinelanguage{JuliaLike}{
  morekeywords={function,for,while,if,else,elseif,end,begin,do,struct,return,where,eachindex,axes,true,false,nothing,Iterators,reverse},
  sensitive=true,
  morecomment=[l]{\#},
  morestring=[b]",
}

\lstset{
  language=JuliaLike,
  basicstyle=\ttfamily\small,
  keywordstyle=\color{codekw}\bfseries,
  commentstyle=\color{codecomment}\itshape,
  backgroundcolor=\color{codebg},
  numbers=left,
  numberstyle=\tiny,
  frame=single,
  breaklines=true,
  showstringspaces=false,
  columns=fullflexible,
  xleftmargin=1.5em,
}

\sloppy
\newcommand{\E}{\mathbb{E}}
\newcommand{\Var}{\mathrm{Var}}
\newcommand{\ESS}{\mathrm{ESS}}

\title{From Code to Mathematics: \\ The Unified Particle Filter in \texttt{pfilter.jl}}
\author{Debsurya}
\date{September 8, 2026}

\begin{document}
\maketitle

\begin{abstract}
This document gives a line-by-line mathematical reading of the unified particle
filter implemented in \texttt{src/pfilter.jl} of the \texttt{PartiallyObservedMarkovProcesses.jl}
package. The filter generalizes the classical bootstrap particle filter along two
axes: an effective-sample-size \emph{trigger} that governs whether resampling occurs
at a given observation time, and a tempering parameter \emph{target} ($\beta$) that
governs how much of a particle's weight is discharged into selection probability
versus retained on the surviving particle. We derive the algorithm from first
principles and show that the two extreme settings recover the classical
bootstrap filter and the fully weighted (never-resampling) filter
respectively. Three genuine defects turned up in the course of this
derivation and were corrected in the code. An earlier translation of the
tempered resampling step indexed the retained weight by selection slot
rather than by selected ancestor (Section~\ref{sec:resample}), and the
terminal ancestry draw used to report a simulated trajectory sampled
uniformly over terminal particle indices rather than in proportion to their
weight (Section~\ref{sec:trace}). The third, and the principal subject of
this document, is that power-tempered resampling discards a normalizing
constant $C$ whenever $0<\beta<1$: Section~\ref{sec:credit} derives $C$
exactly, from the same importance-sampling argument used for the other two
defects, shows by an exact two-particle calculation that dropping it biases
the running log likelihood low by $\Theta(1/n)$ per resampling step, and
confirms the bias with an end-to-end filtering test built specifically to
be sensitive to it. This finding was itself derived once before, implemented,
and then wrongly retracted on the strength of an underpowered numerical
check; Section~\ref{sec:history} narrates that episode and draws the general
lesson about what a null Monte Carlo comparison can and cannot show.
Section~\ref{sec:kalman} gives the general validation of the filter against
the exact Kalman likelihood for the Gompertz model.
\end{abstract}

\tableofcontents

\section{The POMP model and the bootstrap particle filter}

\subsection{The partially observed Markov process}

A partially observed Markov process (POMP) consists of a latent Markov process
$\{X_t\}$ together with a sequence of noisy observations $\{Y_n\}_{n=1}^N$ made
at times $t_1 < t_2 < \cdots < t_N$. Write $x_{0:N} = (x_0,x_1,\dots,x_N)$ for a
realization of the latent process, sampled at $t_0 < t_1 < \cdots < t_N$, and
$y_{1:N}=(y_1,\dots,y_N)$ for the data. The model is specified by three
components:

\begin{itemize}
\item an \emph{initial density}, from which $X_0$ is drawn: $X_0 \sim r_{\mathrm{init}}(\cdot;\theta)$;
\item a \emph{process transition density} $f(x_n \mid x_{n-1};\theta)$, giving the
Markov transition law from $t_{n-1}$ to $t_n$;
\item a \emph{measurement density} $g(y_n \mid x_n;\theta)$, giving the conditional
law of the observation at $t_n$ given the latent state there, with observations
conditionally independent given the latent path.
\end{itemize}

The likelihood of the data under parameter $\theta$ is the marginal density
\[
\mathcal{L}(\theta) = \int \cdots \int r_{\mathrm{init}}(x_0;\theta)
\prod_{n=1}^N f(x_n\mid x_{n-1};\theta)\, g(y_n\mid x_n;\theta) \; dx_0\, dx_1 \cdots dx_N,
\]
an integral of dimension proportional to $N$ that is analytically intractable for
essentially every model of interest beyond the linear-Gaussian case. In the code,
$r_{\mathrm{init}}$, the process model, and $g$ are supplied as the basic
components \verb|rinit|, \verb|rprocess|, and \verb|logdmeasure|; the package
requires only that one can \emph{simulate} from $r_{\mathrm{init}}$ and from
$f(\cdot\mid x_{n-1};\theta)$, and \emph{evaluate} $g(y_n\mid x_n;\theta)$
pointwise — the plug-and-play property that motivates the sequential Monte Carlo
approach below.

\subsection{The bootstrap particle filter}

The bootstrap particle filter (Gordon, Salmond \& Smith, 1993) estimates
$\mathcal{L}(\theta)$ by propagating a swarm of $J$ particles
$\{X_n^{1:J}\}$ through the filtering recursion, using the process model itself
as the proposal (importance) distribution. Initialize
$X_0^{j} \sim r_{\mathrm{init}}(\cdot;\theta)$, $j=1,\dots,J$, independently.
At each observation time $n=1,\dots,N$:

\begin{enumerate}
\item \textbf{Prediction.} Propagate each particle forward through the process
model: $X_n^{j} \sim f(\cdot \mid X_{n-1}^{j};\theta)$, independently across $j$.
\item \textbf{Weighting.} Assign each predicted particle an importance weight
equal to the measurement density evaluated at the data:
\[
w_n^{j} = g(y_n \mid X_n^{j};\theta).
\]
\item \textbf{Resampling.} Draw $J$ indices $\xi_n^1,\dots,\xi_n^J$ with
$\Pr(\xi_n^k = j) \propto w_n^j$, and set the filtered particle
$X_n^{k} \leftarrow X_n^{\xi_n^k}$ for each $k$, discarding the weights (each
surviving particle now carries equal weight $1/J$).
\end{enumerate}

Because the process model is used as the proposal, the importance weight is
exactly the measurement density — no Radon-Nikodym correction between proposal
and target transition kernels is needed. The conditional likelihood contributed
at time $n$ (the density of $y_n$ given $y_{1:n-1}$) is estimated, unbiasedly,
by the empirical mean of the weights:
\[
\widehat{\mathcal{L}}_n = \frac{1}{J}\sum_{j=1}^J w_n^j,
\]
and, since $\mathcal{L}(\theta) = \prod_{n=1}^N \mathcal{L}_n(\theta)$ by the
Markov property, the overall likelihood is estimated by the product
\[
\widehat{\mathcal{L}}(\theta) = \prod_{n=1}^N \widehat{\mathcal{L}}_n
= \prod_{n=1}^N \left( \frac{1}{J}\sum_{j=1}^J w_n^j \right).
\]
This estimator is unbiased for $\mathcal{L}(\theta)$ (Del Moral, 2004) — a
non-trivial fact, since it is a product of ratios of dependent random
quantities, not simply an average of independent terms. The remainder of this
document is concerned with a single question: which modifications of steps~1–3
above preserve this unbiasedness property, and what bookkeeping does the
implementation require to make sure they do.

The package works throughout on the log scale, both for numerical stability
(weights routinely span many orders of magnitude) and because the overall log
likelihood decomposes as a sum,
\[
\log \widehat{\mathcal{L}}(\theta) = \sum_{n=1}^N \log \widehat{\mathcal{L}}_n,
\]
of the \emph{conditional log likelihoods} — exactly the quantity accumulated,
observation time by observation time, in the array \texttt{cond\_logLik} in
\texttt{pfilter.jl}.

\section{Power-tempered systematic resampling}
\label{sec:resample}

\subsection{Tempering the resampling step}

The classical resampling step (step~3 above) discharges the entire weight
$w^j$ of a particle into its selection probability. The unified filter instead
factors each weight as
\[
w = w^{1-\beta} \cdot w^{\beta}, \qquad \beta \in [0,1],
\]
and uses only the first factor to drive selection: particle $j$ is selected
with probability proportional to $(w^j)^{1-\beta}$, and each selected particle
\emph{retains}, rather than discards, the second factor $(w^j)^{\beta}$ as its
new weight — renormalized, as described below, to have unit mean over the
particle swarm. The parameter $\beta$ is called \texttt{target} in the code
(and denoted $\beta$ throughout this document to keep the two roles — the
resampling parameter and the code identifier for it — visually distinct).

\begin{itemize}
\item At $\beta = 0$: selection probability $\propto w^1 = w$, retained weight
$\propto w^0 = 1$. This is exactly the classical bootstrap resampling step:
particles are chosen proportional to their full weight and emerge with equal
weight.
\item At $\beta = 1$: selection probability $\propto w^0 = 1$, retained weight
$\propto w^1 = w$. Selection is now \emph{uninformative} — every particle is
equally likely to be chosen regardless of its weight — while the full weight is
carried forward unchanged on the selected copy. No information is discarded;
this is resampling in name only, and is used to reshuffle particle identity
(for bookkeeping of ancestry) without touching the weighted representation.
\item At intermediate $\beta$: a genuine interpolation. Some of the weight's
information is spent reducing the variance of the particle swarm's spatial
distribution (via selection), and the remainder is preserved as an informative
weight on the surviving particles, to be combined multiplicatively with weights
accrued at future observation times.
\end{itemize}

\subsection{The four-argument \texttt{systematic\_resample!}}

The tempered resampling step is implemented in \texttt{src/pfilter.jl},
lines 457--491:

\begin{lstlisting}[caption={\texttt{systematic\_resample!}, power-tempered systematic resampling (\texttt{pfilter.jl}, lines 457--491)}]
systematic_resample!(
    p::AbstractArray{I,1},
    w::AbstractArray{W,1},
    ucum::AbstractArray{W,1},
    beta::W,
) where {I,W} = begin
    @assert length(ucum)==length(w)==length(p)
    s::W = 0
    alpha = 1-beta
    @inbounds for j in eachindex(w)
        s += w[j]^alpha
        ucum[j] = s
    end
    n::I = length(w)
    i::I = 1
    du::W = s/n
    u::W = -du*rand(W)
    @inbounds for j in eachindex(p)
        u += du
        while (u > ucum[i] && i < n)
            i += 1
        end
        p[j] = i
    end
    stot::W = s
    @inbounds for j in eachindex(p)
        ucum[j] = w[p[j]]^beta
    end
    @inbounds for j in eachindex(w)
        w[j] = ucum[j]
    end
    m::W = mean(w)
    w ./= m # subsequent steps rely on the weights having unit mean
    log(m*stot/n)
end
\end{lstlisting}

We map this line by line onto the mathematics of systematic resampling
(Kitagawa, 1996), the standard low-variance alternative to multinomial
resampling.

\paragraph{Cumulative sum of tempered weights (lines 466--469).} Writing
$w^{1:J}$ for the incoming (linear, not necessarily unit-mean) weight vector
and $\alpha = 1-\beta$, the code computes the cumulative sums
\[
S'_j = \sum_{i=1}^{j} (w^i)^{\alpha}, \qquad j=1,\dots,J, \qquad S' := S'_J = \sum_{i=1}^J (w^i)^{\alpha}.
\]
\texttt{ucum[j]} holds $S'_j$; \texttt{s} accumulates to $S'$ at the end of the
loop.

\paragraph{The systematic sweep (lines 470--480).} A single uniform random
offset $u_0 \sim \mathrm{Unif}(0,S'/J)$ is drawn (line 473, as
\texttt{-du*rand(W)}, so that the first increment lands at exactly this
offset), and $J$ equally spaced points
\[
u_k = u_0 + k\cdot \frac{S'}{J}, \qquad k=1,\dots,J,
\]
are swept against the cumulative sum $S'_{1:J}$: the ancestor index for
selection $k$ is
\[
p_k = \min\{\, j : S'_j \ge u_k \,\}.
\]
This is systematic resampling on the tempered weights $(w^j)^{1-\beta}$: each
ancestor $j$ is selected with multiplicity concentrated near
$J\,(w^j)^{1-\beta}/S'$, and, unlike independent multinomial draws, the
selection has minimal variance subject to these marginal selection
probabilities — a single sorted sweep rather than $J$ independent draws.
\texttt{p[j]} on exit is the array $\xi^{1:J} = (p_1,\dots,p_J)$ of selected
ancestor indices; \texttt{stot} (line 481) holds $S'$ itself, carried
forward for use in the normalizing constant computed at the end of the
routine (Section~\ref{sec:credit}).

\paragraph{Ancestor-indexed retained weight (lines 482--484).} This is the
step at the mathematical heart of the tempered scheme, and the one most easily
gotten wrong (Section~\ref{sec:trace} below describes an analogous error in
ancestry tracing). Each original particle's weight factored as
$w^j = (w^j)^{1-\beta}(w^j)^{\beta}$; the first factor has already been fully
discharged into the selection step above (a particle with a large
$(w^j)^{1-\beta}$ is selected many times). The remaining factor, $(w^j)^\beta$,
is the piece of information not yet used, and it belongs to the \emph{particle},
not to the \emph{selection slot} it lands in. Consequently the retained weight
of the $k$-th selected particle must be indexed by its ancestor:
\[
\tilde{w}_k = (w^{p_k})^{\beta}, \qquad k=1,\dots,J,
\]
exactly what line 483 computes:
\texttt{ucum[j] = w[p[j]]\textasciicircum{}beta}. (\texttt{ucum} is reused here
purely as scratch space, unrelated to its earlier role as the cumulative-sum
array.) The alternative — raising the weight vector to the power $\beta$
\emph{before} indexing by \texttt{p}, i.e. computing
$w[j] \leftarrow w[j]^\beta$ in place and then permuting — assigns to
selection slot $k$ the quantity $(w^k)^\beta$, the retained-weight power of
whichever particle originally occupied position $k$, which is in general
unrelated to the particle that was actually selected there. This positional
form was the form found in an earlier translation of the algorithm, at Aaron
King's commit \texttt{1610d34} (see the memo in this directory, item 2), and
has since been fixed there independently, at commit \texttt{1d8dbc0}, in a
form equivalent to the one below. It targets the wrong measure whenever
$\beta>0$, as Section~\ref{sec:trace} demonstrates concretely by a
two-particle example in the closely analogous setting of the terminal
ancestry draw.

\paragraph{Writing back and renormalizing (lines 485--489).} The retained
weights computed in the preceding loop are written into \texttt{w}
(lines 485--487), and \texttt{w} is then divided by its own realized mean
\[
m = \frac1J\sum_k \tilde w_k,
\]
computed at line 488 (\texttt{m::W = mean(w)}) and applied at line 489
(\texttt{w ./= m}), so that on exit \texttt{w} has unit mean over the $J$
particles. This is an implementation convention, not a mathematical
necessity: the routine could equally well leave the retained weights
unnormalized. The choice to renormalize is made because \emph{every other
stage of the filter} — in particular \texttt{compute\_ess\_logLik!},
discussed in Section~\ref{sec:ess} — assumes an entering weight vector with
unit mean, so that the log of its arithmetic mean, needed for the
conditional log likelihood, can be read off directly. Keeping this invariant
globally is what makes it possible for the filter to treat "just resampled"
and "several steps since the last resampling" states identically at every
subsequent step.

\paragraph{The normalizing constant, returned (line 490).} Renormalizing to
unit mean divides every retained weight — and with it the swarm's total
mass — by the common factor $m$. Because the sampling measure driving
selection was $q_j\propto (w^j)^{1-\beta}$ rather than the weights
themselves, this renormalization removes real likelihood-scale information,
not merely a bookkeeping convenience: Section~\ref{sec:credit} shows that
the quantity
\[
C = \frac{m\cdot S'}{J}
\]
is exactly the resampling step's own contribution to the normalizing
constant of the filter's target measure, that it has conditional mean one
but is not independent of what it multiplies, and that dropping it biases
the running log likelihood. The routine accordingly returns $\log C =
\log(m\cdot\texttt{stot}/n)$ (line 490, with \texttt{n} the particle count
$J$) rather than \texttt{nothing}, and the caller
(Section~\ref{sec:step}, line 315) adds this return value into the running
log likelihood.

\section{The normalizing constant discarded at partial resampling}
\label{sec:credit}

\subsection{The telescoping identity}

Consider two consecutive observation times, with the filter having just
computed, at time $n-1$, a properly weighted particle representation of the
filtering distribution: a swarm $\{(X_{n-1}^j, \omega_{n-1}^j)\}_{j=1}^J$ with
$\E\left[\frac1J\sum_j \omega_{n-1}^j h(X_{n-1}^j)\right] = \E[h(X_{n-1})\mid
y_{1:n-1}]$ for every test function $h$, and normalized so
$\frac1J\sum_j \omega_{n-1}^j = 1$ (unit mean). Propagating each particle
through the process model and weighting by the measurement density at time
$n$ produces an updated properly weighted swarm for the joint smoothing
distribution of $X_n$ given $y_{1:n}$, with un-normalized weights
$\omega_{n-1}^j \cdot g(y_n\mid X_n^j;\theta)$, and — because the incoming
swarm already had unit mean and represents the filtering measure correctly —
the conditional likelihood contribution is exactly the mean of these updated
weights:
\[
\mathcal{L}_n \;=\; \frac1J \sum_{j=1}^J \omega_{n-1}^j\, g(y_n\mid X_n^j;\theta).
\]
This is the content of the telescoping (or "sequential") identity
$\mathcal{L}(\theta)=\prod_n \mathcal{L}_n(\theta)$ specialized to a
weighted particle representation: \emph{as long as the entering weights have
unit mean and properly represent the filtering measure, the conditional
likelihood at each step is recovered simply by averaging the updated,
un-normalized weights} — no correction term is needed for the propagate/weight
half of the recursion. The delicate step is what happens next, at
resampling: the identity survives only if the swarm handed to the next
observation time remains a properly weighted representation of the same
measure, up to a normalizing constant that must be tracked, not discarded.
The rest of this section derives that constant exactly and shows what goes
wrong when it is dropped.

\subsection{Proper weighting: a sufficient condition}
\label{sec:proper-weighting}

After weighting, resampling replaces the properly weighted swarm
$\{(X_n^j, w_n^j)\}$ — where, dropping the accumulated-weight notation
$\omega$ used above, $w_n^j$ denotes the current, possibly un-normalized
weight at time $n$ — with a new swarm $\{(X_n^{p_k}, \tilde w_{n}^k)\}_{k=1}^J$
obtained via power-tempered systematic resampling with parameter $\beta$. One
sufficient condition for the telescoping identity to continue to hold
undisturbed at the next step is that the new swarm remain \emph{properly
weighted} for the same filtering measure in the classical importance-sampling
sense: for every test function $h$,
\[
\E\left[\frac1J \sum_{k=1}^J \tilde w_n^k\, h(X_n^{p_k}) \;\middle|\; X_n^{1:J}, w_n^{1:J}\right]
= \frac{1}{S}\sum_{j=1}^J w_n^j\, h(X_n^j), \qquad S := \sum_{j=1}^J w_n^j,
\]
i.e. the resampled swarm remains an unbiased (times a known constant)
representation of the same weighted empirical measure. A direct calculation —
using $\Pr(p_k = j) = (w_n^j)^{1-\beta}/S'$ with $S' = \sum_i (w_n^i)^{1-\beta}$,
and retained weight $\tilde w_n^k = (w_n^{p_k})^\beta / m$ where $m$ is
whatever normalizing constant is applied — shows
\[
\E\left[\tilde w_n^k\, h(X_n^{p_k}) \mid X_n^{1:J}, w_n^{1:J}\right]
= \frac{1}{m}\sum_{j=1}^J \frac{(w_n^j)^{1-\beta}}{S'}\,(w_n^j)^\beta\, h(X_n^j)
= \frac{1}{m\,S'} \sum_{j=1}^J w_n^j\, h(X_n^j).
\]
Comparing with the target above, this sufficient condition holds exactly when
$m S' = S$, i.e. when the retained weight of the $k$-th selected particle is
taken to be
\[
\tilde w_n^k = (w_n^{p_k})^\beta \cdot \frac{S'}{S},
\]
which, when the incoming weights already have unit mean ($S=J$), is
$(w_n^{p_k})^\beta\cdot S'/J$ — this is the correct, and the only, assignment
that makes the resampled swarm properly weighted in the sense above. It is
not, however, the assignment the code makes, and it is worth being precise
about the gap between the two before asking whether that gap actually
matters.

\subsection{The normalizing constant, exactly}
\label{sec:removed-factor}

Rather than the properly weighted assignment $\tilde w_n^k = (w_n^{p_k})^\beta
\cdot S'/S$ of Section~\ref{sec:proper-weighting}, the code renormalizes to
\emph{unit mean}, dividing every retained weight $(w_n^{p_k})^\beta$ by the
realized sample mean
\[
m = \frac1J\sum_{k=1}^J (w_n^{p_k})^\beta.
\]
Both are valid ways to fix the overall scale of the resampled swarm, and
they generally disagree: $S/S'$ is a deterministic function of the
pre-resampling weights alone, while $m$ is a realized average over the $J$
systematic selections. Writing $R_k = w_n^{p_k}\big/(J\,q_{p_k})$ with
$q_j=(w_n^j)^{1-\beta}/S'$ for the importance weight the proper-weighting
argument assigns to selection $k$ — so that $R_k = (S'/J)\,(w_n^{p_k})^\beta$
— the code's realized weight is $V_k = (w_n^{p_k})^\beta/m$, related to the
properly weighted $R_k$ by
\[
V_k = \frac{R_k}{C}, \qquad C := \frac{m\,S'}{J} \quad \text{(using $S=J$
under the incoming unit-mean convention).}
\]
$C$ is exactly the resampling step's own contribution to the normalizing
constant of the unnormalized (Feynman--Kac) measure the filter is
estimating: it is the sample mean of the properly weighted quantities $R_k$
themselves,
\[
C = \frac1J\sum_{k=1}^J R_k,
\]
as a short calculation confirms: $\frac1J\sum_k R_k = \frac1J\sum_k
(S'/J)(w_n^{p_k})^\beta = (S'/J)\cdot m$, matching the definition above. This
is not a bookkeeping curiosity — $C$ is a piece of the likelihood, on exactly
the same footing as the mean weight computed at every other observation
time in \texttt{compute\_ess\_logLik!} (Section~\ref{sec:ess}) — and
renormalizing $V_k=R_k/C$ to unit mean, as the code does, silently divides
it out. Proposition~\ref{prop:credit} below records exactly when $C$ is
forced to $1$ and when it is not.

\begin{proposition}
\label{prop:credit}
Suppose the incoming weight vector $w^{1:J}$ has unit mean,
$\frac1J\sum_j w^j = 1$ (as guaranteed by \texttt{compute\_ess\_logLik!},
Section~\ref{sec:ess}, on entry to resampling). Let $S' = \sum_j
(w^j)^{1-\beta}$ and let $m$ be the realized mean, over the $J$ systematic
selections $p_{1:J}$, of $(w^{p_k})^\beta$, and $C = mS'/J$. Then:
\begin{enumerate}[label=(\alph*)]
\item $C = 1$ identically (for every realization of the resampling
randomness) when $\beta = 0$;
\item $C = 1$ identically when $\beta = 1$;
\item $\E[C\mid w^{1:J}] = 1$ exactly, for every $\beta\in[0,1]$ and every
realization of the pre-resampling weights; but for $0 < \beta < 1$, $C$ is
generically a nondegenerate random variable, not almost-surely equal to
$1$.
\end{enumerate}
\end{proposition}

\begin{proof}
(a) At $\beta=0$: $S' = \sum_j (w^j)^1 = \sum_j w^j = J$ (unit mean), and
$m = \frac1J\sum_k (w^{p_k})^0 = \frac1J\sum_k 1 = 1$. So $C = 1\cdot J/J
= 1$, for every realization — this is simply the statement that classical
resampling to equal weights removes no likelihood information: the entire
mass $S=J$ is preserved exactly, and no renormalization is needed ($m=1$
trivially).

(b) At $\beta=1$: $S' = \sum_j (w^j)^0 = J$, and $m = \frac1J\sum_k
(w^{p_k})^1 = \frac1J\sum_k w^{p_k}$. Since selection at $\beta=1$ is uniform
systematic resampling (every $(w^j)^0=1$, so all ancestors are equally
likely), and systematic resampling from a uniform base measure over $J$
categories, each with population fraction exactly $1/J$, reproduces each
index exactly once (systematic resampling against a perfectly uniform
cumulative distribution is a bijection, not merely unbiased): $p$ is a
permutation of $1,\dots,J$, hence $m = \frac1J \sum_k w^{p_k} = \frac1J\sum_j
w^j = 1$ (unit mean, again). So $C = 1\cdot J/J = 1$, for every realization:
at $\beta=1$ resampling is a relabeling that touches neither the weights nor
their sum.

(c) $S'$ is a deterministic function of the entering weights, so the
resampling randomness enters only through $m$. Writing $q_j = (w^j)^{1-\beta}/S'$
for the selection probabilities, $\E[m\mid w^{1:J}] = \sum_j q_j (w^j)^\beta
= \frac1{S'}\sum_j w^j = J/S'$ (using $\sum_j w^j = J$), so $\E[C\mid
w^{1:J}] = (J/S')\cdot(S'/J) = 1$, exactly, for every $\beta$ — this
computation makes no use of $\beta\in\{0,1\}$ and so subsumes (a) and (b) as
the special, degenerate cases in which $m$ (and hence $C$) is not merely
mean-one but almost surely equal to $1$. For $0<\beta<1$, by contrast, the
selection multiplicities are not uniform and $m$ varies with the resampling
draw: $C$ is a genuine random variable with mean $1$, not the constant $1$.
\end{proof}

\begin{corollary}
\label{cor:extremes}
$C$ of Proposition~\ref{prop:credit} is identically $1$ whenever the filter
is run at its classical default ($\beta=$~\texttt{target}~$=0$) or at the
fully weighted extreme ($\beta=1$); at those two settings, restoring $C$ or
dropping it makes no difference of any kind, to any realization —
multiplying by exactly $1$ changes nothing. Only at intermediate $\beta$
does $C$ ever differ from $1$, and it is precisely there, as the next
subsection shows, that dropping it costs more than the conditional-mean-one
property of part (c) might suggest.
\end{corollary}

\subsection{Conditional mean one is not enough}
\label{sec:telescoping-route}

The code, as it stands, does not restore $C$: \texttt{systematic\_resample!}
returned \texttt{nothing}, and $C$ was never multiplied into the likelihood.
This was precisely the bookkeeping already used, for entirely unrelated
reasons of implementation convenience, in the R \texttt{pomp} package's
particle filter (\texttt{wpfilter.c}): weights are carried forward without
restoring $C$ at resampling, and each observation time's conditional log
likelihood is instead computed as the log of a \emph{ratio} of successive
total masses, \texttt{loglik = maxlogw + log(w) - log(oldw)}
(\texttt{wpfilter.c}, line 110), where \texttt{oldw} is simply whatever
total mass the carried weights happen to have on entry to that step —
resampled and renormalized, or not. $C$ is dropped, not restored, by this
route too: the mass change across a resampling step is never entered into
the likelihood at all, and the next observation's increment is taken
relative to whatever baseline the carried weights land on. Aaron King's R
implementation and, until now, both Julia translations compute, and always
have computed, the same estimator, and it is biased.

Proposition~\ref{prop:credit}(c) shows $\E[C\mid w^{1:J}]=1$ for every
$\beta$. It is tempting to conclude from this alone that dropping $C$ costs
only variance: a factor with conditional expectation one, multiplied into
an otherwise-unbiased estimator, cannot change that estimator's expectation
— \emph{if} the factor is independent of everything else the estimator
depends on. $C$ is not independent of what it multiplies. $C = m S'/J$ is a
function of the realized selections $p_{1:J}$, i.e.\ of \emph{which
ancestors survive resampling}, and those same ancestors are exactly the
particles that generate every future contribution to the likelihood: the
process draws at every subsequent observation time, and the measurement
densities evaluated against them, all descend from $X_n^{p_{1:J}}$. Knowing
$\E[C]=1$ marginally, or even conditionally on the pre-resampling weights,
does not give $\E[C\cdot\ell]=\E[\ell]$ for a downstream quantity $\ell$
that is correlated with $C$ through the shared ancestry — and a resampling
step that determines which particles survive is, by construction,
correlated with everything those particles subsequently produce. Dropping
$C$ therefore biases the likelihood estimate, by $\Theta(1/J)$ per
resampling step at $0<\beta<1$; Corollary~\ref{cor:extremes} already
confines the effect to that range, since $C\equiv1$ pins the bias to
exactly zero at $\beta\in\{0,1\}$. Section~\ref{sec:numerical-credit} makes
this concrete with an exact calculation small enough to enumerate by hand,
and then with an end-to-end filtering test.

\subsection{An exact enumeration, and an end-to-end test}
\label{sec:numerical-credit}

\paragraph{A two-particle enumeration.} Take $J=2$ particles with unit-mean
weights $w=(1.8,0.2)$ and $\beta=1/2$. Then $S'=\sqrt{1.8}+\sqrt{0.2}\approx
1.7889$, $q=(0.75,0.25)$, and two-particle systematic resampling draws the
ancestry $(1,1)$ with probability $1/2$ and $(1,2)$ with probability $1/2$
(the two possible outcomes of a single systematic sweep against $J=2$
categories with these selection probabilities). Under $(1,1)$: $m =
\frac12\big(1.8^{1/2}+1.8^{1/2}\big)=\sqrt{1.8}\approx1.3416$ and
$C=mS'/2\approx1.2$. Under $(1,2)$: $m=\frac12(1.8^{1/2}+0.2^{1/2})
=S'/2\approx0.8944$ and $C\approx0.8$. (Both values of $C$ round to exactly
$1.2$ and $0.8$; this is not a coincidence of rounding but exact arithmetic,
since $\sqrt{1.8}\cdot S'/2$ and $(S'/2)\cdot S'/2$ simplify exactly using
$S'=\sqrt{1.8}+\sqrt{0.2}$ and $\sqrt{1.8}\sqrt{0.2}=0.6$.) Let the future
contribution from the two particles be $g=(1,0)$; the correctly weighted
predictive quantity is $\frac12\sum_i w_i g_i = 0.9$. Dropping $C$ gives
$\frac12(1)+\frac12(0.75)=0.875$; restoring it gives
$\frac12(1.2\cdot1)+\frac12(0.8\cdot0.75)=0.9$, exactly — this has been
verified numerically to six decimal places over $4{,}000{,}000$ replicates
of the resampling step.

\paragraph{An end-to-end test.} \texttt{test/iid.jl} contains two test sets
built around this distinction. The first uses a latent process with no
memory ($X_t$ drawn afresh, independently of $X_{t-1}$, at every step), and
documents that such a model is \emph{blind} to the effect: the one-step
predictive density is then the same for every particle regardless of which
ancestors resampling selects, so the tempered filter is exactly unbiased
whatever it does with $C$, at any number of replicates. The second test
constructs the case the first cannot: a frozen binary latent state,
$X\sim\mathrm{Bernoulli}(1/2)$ with $X_t=X$ for all $t$, so the surviving
particle after resampling determines the entire remainder of the path.
Measurement potentials $g_1(1)=1.8,\,g_1(0)=0.2$ and $g_2(1)=1.0,\,g_2(0)=0.0$
give an exact likelihood $Z=\E[g_1(X)g_2(X)]=(1/2)(1.8)(1.0)=0.9$. At $J=2$
and $\beta=1/2$ — the setting of the hand calculation above, where the exact
predicted value of the dropped-$C$ estimator is $0.8875$ — the test measures
$\E[\hat Z]\approx0.8879$ over $200{,}000$ replicates, a discrepancy of
about $27$ Monte Carlo standard errors from the exact $Z=0.9$; restoring $C$
removes the discrepancy. The same test reports the gap shrinking as $J$
grows, consistent with the $\Theta(1/J)$ scaling of
Proposition~\ref{prop:credit}(c).

\subsection{The code}

\texttt{systematic\_resample!} now returns $\log C$ rather than
\texttt{nothing}, and the caller accumulates it:

\begin{lstlisting}[caption={The end of \texttt{systematic\_resample!} (\texttt{pfilter.jl}, lines 481--490): $C$ computed and returned}]
    stot::W = s
    @inbounds for j in eachindex(p)
        ucum[j] = w[p[j]]^beta
    end
    @inbounds for j in eachindex(w)
        w[j] = ucum[j]
    end
    m::W = mean(w)
    w ./= m # subsequent steps rely on the weights having unit mean
    log(m*stot/n)
end
\end{lstlisting}

\begin{lstlisting}[caption={The caller, \texttt{pfilt\_step\_comps!} (\texttt{pfilter.jl}, line 315): the return value is accumulated}]
        logLik[] += systematic_resample!(p, w, work, target)
\end{lstlisting}

With this bookkeeping, \texttt{cond\_logLik[k]} — no longer left untouched
by resampling — restores the identity
$\widehat{\mathcal{L}}(\theta)=\prod_n \widehat{\mathcal{L}}_n$ as an
unbiased estimator of $\mathcal{L}(\theta)$ uniformly over the trigger/target
parameter space, by the proper-weighting route of
Section~\ref{sec:proper-weighting}: the swarm handed to the next observation
time is, once $C$ is restored, exactly the properly weighted swarm that
section derives, not merely a mean-one rescaling of it. Aaron King's Julia
translation and R \texttt{pomp}'s \texttt{wpfilter.c} still drop $C$;
Section~\ref{sec:kalman} validates the filter as a whole numerically against
an exact reference likelihood, and Section~\ref{sec:history} narrates how
this correction came to be made twice.

\section{ESS-triggered resampling}
\label{sec:ess}

\subsection{Effective sample size}

The effective sample size of a weighted particle swarm $w^{1:J}$ is the
classical diagnostic
\[
\ESS(w^{1:J}) = \frac{\left(\sum_{j=1}^J w^j\right)^2}{\sum_{j=1}^J (w^j)^2},
\]
a quantity ranging from $1$ (all mass on one particle — maximal weight
degeneracy) to $J$ (all weights equal — no degeneracy), invariant to a
uniform rescaling of the weights. It is the standard proxy for "how many
independent, equally weighted particles is this weighted swarm worth,"
and is the basis for adaptive resampling schedules: resample only when the
swarm has become sufficiently degenerate to threaten the accuracy of the
filtering approximation, rather than at every step regardless of need.

\subsection{\texttt{compute\_ess\_logLik!}: the two methods}

Two methods of \texttt{compute\_ess\_logLik!} appear in the code
(lines 330--403), corresponding to the classic and general engines
respectively.

\paragraph{Classic method (lines 330--360), no carried weight.} Given only
the incremental log weights $\log w_n^j = \log g(y_n\mid X_n^j;\theta)$ (no
prior weight carried in, since the classic engine always resamples to equal
weights at the previous step), the routine:
\begin{itemize}
\item finds $\ell_{\max} = \max_j \log w_n^j$ (line 335) and subtracts it from
every log weight (line 344), the standard log-sum-exp stabilization, so that
the largest stabilized weight is $1$ before exponentiating;
\item accumulates $s=\sum_j v^j$ and $ss=\sum_j (v^j)^2$ with $v^j =
\exp(\log w_n^j - \ell_{\max})$ (lines 341--348);
\item sets $\ESS = s^2/ss$ (line 350), exactly the definition above, computed
on the stabilized (equivalently, unstabilized — ESS is scale-invariant)
weights;
\item computes the conditional log likelihood as
$\log\widehat{\mathcal L}_n = \ell_{\max} + \log(s/J)$ (lines 349--352), the
log-sum-exp identity
$\log\!\left(\frac1J\sum_j w_n^j\right) = \ell_{\max}+\log\!\left(\frac1J\sum_j
e^{\log w_n^j - \ell_{\max}}\right)$;
\item overwrites \texttt{logw} with the log weights renormalized to have
mean-one linear scale, $\log w_n^j - \log(s/J)$ (line 353), so that
\emph{even along the path where resampling does not subsequently occur}, the
weights handed onward are on a canonical, likelihood-consistent scale (used,
for instance, directly by \texttt{systematic\_resample!}'s classic sibling
method, or, in the general engine, becoming the next step's carried weight).
\end{itemize}
If $\ell_{\max}$ is not finite (line 340's \texttt{else} branch), every
particle has zero measurement density — the swarm has been exhausted, filter
failure — and the routine reports $\ESS=0$, $\log\widehat{\mathcal L}_n =
-\infty$ (lines 355--357).

\paragraph{General method (lines 368--403), folding in a carried weight.}
The general engine carries a linear, unit-mean weight vector \texttt{w}
forward across observation times whenever resampling does not occur at every
step (Section~\ref{sec:step}); this method folds that carried weight into
the current step's incremental log weight before computing ESS and the
conditional log likelihood, and — critically — writes the folded, renormalized
result \emph{back into} \texttt{w} on exit (line 395,
\texttt{w ./= lik}), maintaining the unit-mean invariant on which every
other routine in the filter relies. Concretely:
\[
\log \tilde w^j = \log w_{\mathrm{carry}}^j + \log w_n^j - \ell_{\max}
\quad\text{(line 384)},
\]
combining the incoming (unit-mean) carried weight multiplicatively with the
new incremental weight, again with the log-sum-exp stabilization
$\ell_{\max}=\max_j \log w_n^j$ (the incremental weight's own maximum, not
the combined one — since the carried weight already has unit mean and no
extreme values by construction). $\ESS$ and $\log\widehat{\mathcal L}_n$ are
computed exactly as before from $s=\sum_j \exp(\log\tilde w^j)$ and
$ss=\sum_j \exp(\log\tilde w^j)^2$, with the comment on line 390 flagging
explicitly that the formula $\log\widehat{\mathcal L}_n = \ell_{\max}+\log(s/J)$
is valid here \emph{only because} the carried weight entered with unit mean —
the same telescoping-identity fact derived in
Section~\ref{sec:credit}. On exit, \texttt{w} is renormalized to unit mean
(line 395) and \texttt{logw} is renormalized to zero mean on the log-linear
scale (line 394), so that whichever code path runs next (resampling or not)
sees a canonically scaled representation.

\subsection{The trigger}

The trigger threshold, \texttt{trigger}~$\in[0,1]$ (denoted here by the same
identifier, since it plays no separate Greek-letter role in the mathematics),
governs whether resampling occurs at all: resampling is invoked whenever
\[
\ESS_n \le \texttt{trigger}\cdot J,
\]
tested at line 314 (\texttt{ess[] \(\le\) trigger*n}), \emph{and} the maximum
log weight is finite (a degenerate/exhausted swarm always triggers the
fallback path, never the resampling path, since there is nothing informative
to resample). \texttt{trigger = 1} resamples unconditionally, since
$\ESS_n\le J$ always; \texttt{trigger = 0} never resamples via this branch
(short of exact degeneracy), since $\ESS_n \ge 1 > 0$ always for a
non-degenerate swarm.

\section{The filtering step}
\label{sec:step}

\texttt{pfilt\_step\_comps!} assembles one full observation-time update:
weighting statistics, the resampling decision, and (when triggered) the
resampling operation itself, dispatched to one of two methods according to
the number of arguments supplied — a compile-time dispatch on whether the
classic or general engine is in use, resolved once, outside the per-step
loop, by the choice made in \texttt{pfilter} at lines 98--136 of which
\texttt{pfilter\_internal!} method to call.

\subsection{Classic path: \texttt{trigger}~$=1$, \texttt{target}~$=0$}

\begin{lstlisting}[caption={Classic filtering step (\texttt{pfilter.jl}, lines 271--293)}]
pfilt_step_comps!(
    logLik::AbstractArray{W,0},
    ess::AbstractArray{W,0},
    logw::AbstractArray{W,1},
    p::AbstractArray{I,1},
    xp::AbstractArray{X,1},
    xf::AbstractArray{X,1},
    work::AbstractArray{W,1},
    resamp::AbstractArray{Bool,0},
    n::Integer = length(logw),
) where {W<:AbstractFloat,I<:Integer,X<:NamedTuple} = begin
    logwmax = compute_ess_logLik!(ess, logLik, logw)
    if isfinite(logwmax)
        systematic_resample!(p, logw, work)
        @inbounds xf .= xp[p]
        resamp[] = true
    else
        p .= collect(eachindex(p))
        xf .= xp
        resamp[] = false
    end
    nothing
end
\end{lstlisting}

This is the algorithm of Section~1.2 exactly: $\log\widehat{\mathcal L}_n$
and $\ESS_n$ are computed from the incremental weights alone (no carried
weight — every step starts from equally weighted particles, since the
classic engine resamples unconditionally whenever the swarm is non-degenerate),
then classical (untempered, $\beta=0$) systematic resampling is applied
whenever the weights are informative ($\log w_{\max}$ finite), replacing
$X_n^{p} $ for the filtered particles $X_n^f$ (line 285,
\texttt{xf .= xp[p]}). In the degenerate case, the identity permutation is
used and the predicted particles are carried through unchanged as the
filtered particles (lines 288--290) — the swarm has failed, but the routine
does not error, allowing the caller to detect and report filter failure
through $\log\widehat{\mathcal L}_n=-\infty$ downstream.

\subsection{General path: weighted, trigger/target}

\begin{lstlisting}[caption={General filtering step (\texttt{pfilter.jl}, lines 299--324)}]
pfilt_step_comps!(
    logLik::AbstractArray{W,0},
    ess::AbstractArray{W,0},
    logw::AbstractArray{W,1},
    p::AbstractArray{I,1},
    xp::AbstractArray{X,1},
    xf::AbstractArray{X,1},
    w::AbstractArray{W,1},
    work::AbstractArray{W,1},
    trigger::W,
    target::W,
    resamp::AbstractArray{Bool,0},
    n::Integer = length(logw),
) where {W<:AbstractFloat,I<:Integer,X<:NamedTuple} = begin
    logwmax = compute_ess_logLik!(ess, logLik, logw, w)
    if isfinite(logwmax) && ess[] <= trigger*n
        logLik[] += systematic_resample!(p, w, work, target)
        @inbounds xf .= xp[p]
        resamp[] = true
    else
        p .= collect(eachindex(p))
        xf .= xp
        resamp[] = false
    end
    nothing
end
\end{lstlisting}

The general step folds the carried weight \texttt{w} into the incremental
weight via \texttt{compute\_ess\_logLik!}'s second method
(Section~\ref{sec:ess}), then resamples only when both the swarm is
non-degenerate \emph{and} the trigger condition
$\ESS_n\le\texttt{trigger}\cdot J$ is met — otherwise the entire weighted
particle representation is carried forward untouched (lines 319--321), with
\texttt{logLik[]} left at the value set by \texttt{compute\_ess\_logLik!}
alone. When resampling does occur, it is tempered resampling at the supplied
$\beta=$~\texttt{target} (line 315); as Section~\ref{sec:credit} works out
in detail, \texttt{logLik[]} \emph{is} updated by this call — the
normalizing constant $C$ that the resampling routine's own unit-mean
renormalization would otherwise remove is computed inside
\texttt{systematic\_resample!} and added into \texttt{logLik[]} at this
line, restoring the properly weighted representation of
Section~\ref{sec:proper-weighting} rather than leaving $C$ to be silently
absorbed downstream.

Comparing the two methods: setting \texttt{trigger}~$=1$ and
\texttt{target}~$=0$ in the general path reduces it, step for step, to the
classic path — $\ESS_n\le 1\cdot J$ always holds, so resampling always
triggers, and $\beta=0$ tempered resampling coincides with classical
resampling exactly, $C$ of Proposition~\ref{prop:credit}(a) being
identically one at $\beta=0$, so no discrepancy of any kind arises between
the two paths at this setting. The classic method is retained separately in
the code as a specialized, slightly cheaper implementation of this special
case (no carried-weight bookkeeping, no trigger test to evaluate), not
because the mathematics differs.

\section{Weighted terminal ancestry draw}
\label{sec:trace}

\subsection{The problem}

At the end of a filtering run, the package reports a single simulated
trajectory $X_{0:N}^\star$ from (an approximation to) the smoothing
distribution, obtained by tracing the ancestry of one randomly selected
terminal particle back through the recorded resampling permutations
\texttt{perm}. When the terminal particle swarm is \emph{equally} weighted —
guaranteed under the classic engine, where every step resamples to equal
weights — a uniform draw over the $J$ terminal particles is exactly correct:
each terminal particle already represents an equal share of the filtering
measure's mass. But whenever \texttt{target}~$>0$ or \texttt{trigger}~$<1$,
the terminal swarm can carry unequal weights $w^{1:J}$ (accumulated exactly
as in Sections~\ref{sec:resample}--\ref{sec:ess}), and a uniform draw over
particle \emph{indices} no longer targets the filtering measure: it
implicitly reweights every particle to $1/J$, discarding the very
information the weights encode.

\subsection{A two-particle counterexample}

\begin{example}
Take $J=2$ terminal particles with weights $w^1 = 3$, $w^2=1$ (unit mean:
$\tfrac12(3+1)=2\ne1$ in this raw form, but the argument is unaffected by an
overall rescaling — take $w^1=1.5,w^2=0.5$ for the unit-mean version if
preferred). The properly weighted representation of the filtering
distribution at this final time is
\[
\pi(\cdot) \;\propto\; w^1\,\delta_{X^1}(\cdot) + w^2\,\delta_{X^2}(\cdot)
\;=\; \tfrac34\,\delta_{X^1}(\cdot) + \tfrac14\,\delta_{X^2}(\cdot)
\]
after normalizing: particle $1$ should be selected as the terminus of the
reported trajectory with probability $3/4$, particle $2$ with probability
$1/4$. A uniform draw, $j\sim\mathrm{Unif}\{1,2\}$, instead selects each
particle with probability $1/2$ — targeting the \emph{unweighted empirical
measure} $\tfrac12\delta_{X^1}+\tfrac12\delta_{X^2}$, a different
distribution whenever $w^1\ne w^2$, and one with no claim to approximating
the filtering distribution of interest. The discrepancy is not a
higher-order or asymptotically vanishing effect: it persists at every $J$,
since it is a first-moment (bias in the selection probability), not a
sampling-variance, error.
\end{example}

This is precisely the defect flagged, as a marginal question rather than a
firm claim, against the uniform draw in an earlier version of
\texttt{trace\_ancestry!} (memo, item 4): "I noticed your marginal note
there, questioning whether the first index in the ancestry is right — I
suspect this is exactly the issue."

\subsection{The corrected, five-argument \texttt{trace\_ancestry!}}

\begin{lstlisting}[caption={Weighted ancestral-lineage tracing (\texttt{pfilter.jl}, lines 508--544)}]
trace_ancestry!(
    traj::AbstractArray{X,1},
    filt::AbstractArray{X,2},
    perm::AbstractArray{I,2},
    w::AbstractArray{W,1},
    work::AbstractArray{W,1},
) where {X,I<:Integer,W<:AbstractFloat} = begin
    @assert size(traj,1)==size(perm,1)
    @assert size(filt)==size(perm)
    @assert length(w)==size(perm,2)
    n = length(w)
    wmax::W = -Inf
    for k in eachindex(w)
        @inbounds wmax = (w[k] > wmax) ? w[k] : wmax
    end
    j::I = if isfinite(wmax)
        s::W = 0
        for k in eachindex(w)
            @inbounds v::W = w[k]
            s += v
            @inbounds work[k] = s
        end
        u::W = s*rand(LogLik)
        jj::I = 1
        while (u > work[jj] && jj < n)
            jj += 1
        end
        jj
    else
        rand(axes(perm,2))
    end
    for i in Iterators.reverse(axes(perm,1))
        @inbounds traj[i] = filt[i,j]
        @inbounds j = perm[i,j]
    end
    j
end
\end{lstlisting}

\paragraph{Cumulative weight construction (lines 523--529).} Given the
terminal (unit-mean, linear) weight vector $w^{1:J}$, the routine forms the
cumulative sums $W_k = \sum_{i=1}^k w^i$, $k=1,\dots,J$ — \texttt{work[k]} —
exactly as \texttt{systematic\_resample!} does for the tempered selection
weights, but here for a \emph{single} draw rather than a systematic sweep of
$J$ draws.

\paragraph{Weighted single draw (lines 530--535).} A draw $u\sim
\mathrm{Unif}(0,W_J)$ is made (line 530, \texttt{s*rand(LogLik)}, with
$W_J=s=\sum_k w^k$), and the smallest index $j$ with $W_j \ge u$ is found by
linear search — an inverse-CDF draw from the discrete distribution with
$\Pr(J=k)=w^k/\sum_i w^i$, exactly the weighted selection probability
required by the properly weighted representation. This is the correction to
the counterexample above: with $w^1=3,w^2=1$, $u\sim\mathrm{Unif}(0,4)$, and
particle $1$ is selected whenever $u\in(0,3]$ — probability $3/4$, as
required.

\paragraph{Degenerate fallback (line 537).} If every weight is
$-\infty$ on the log scale (here checked via a $-\infty$ sentinel on the
linear scale, \texttt{wmax} not finite — mirroring the fallback convention
used throughout \texttt{pfilter.jl}), the routine falls back to a uniform
draw, matching the equal-weight case and the behavior of the unweighted,
four-argument \texttt{trace\_ancestry!} (lines 483--496) used by the classic
engine, where every terminal particle carries equal weight by construction
and this method is unnecessary (and unused).

\paragraph{Backward trace (lines 539--543), identical in both methods.} Once
the terminal index $j$ is fixed (by whichever draw was appropriate), the
trajectory is traced backward through the recorded ancestor permutations
exactly as in the unweighted case: at each observation time $i$, from the
last to the first, $\texttt{traj}[i]\leftarrow \texttt{filt}[i,j]$, then
$j\leftarrow \texttt{perm}[i,j]$, walking the lineage of ancestors back to
time $t_0$. This part of the routine is unaffected by weighting — the
weighting question is entirely about \emph{which} terminal lineage to
initiate, not how to trace one once chosen — and is shared, verbatim in
structure, between the two methods.

\section{The Gompertz--Kalman validation}
\label{sec:kalman}

\subsection{The log-Gompertz linear-Gaussian reduction}

The Gompertz population model, in the parametrization used by the package's
built-in example, is
\[
X_t = X_{t-1}^{S}\, K^{1-S}\, \varepsilon_t, \qquad S = e^{-r}, \qquad
\varepsilon_t \sim \mathrm{LogNormal}(0,\sigma_p^2),
\]
with a lognormal measurement model, $\mathrm{pop}_t \mid X_t \sim
\mathrm{LogNormal}(\log X_t, \sigma_m^2)$. Writing $Z_t = \log X_t$, taking
logs of the process equation gives an exactly linear, Gaussian
autoregression:
\[
Z_t = (1-S)\log K + S\,Z_{t-1} + w_t, \qquad w_t \sim \mathrm{Normal}(0,\sigma_p^2),
\]
and the measurement model becomes, on the log scale, another linear-Gaussian
observation equation:
\[
\log(\mathrm{pop}_t) = Z_t + v_t, \qquad v_t \sim \mathrm{Normal}(0,\sigma_m^2).
\]
This is a scalar linear-Gaussian state-space model, for which the likelihood
of $\log(\mathrm{pop}_{1:N})$ is available exactly, with no Monte Carlo
error, via the Kalman filter — a rare and valuable property for validating a
Monte Carlo filter implementation, since the comparison is against a
deterministic reference rather than a second noisy estimate.

In the package's built-in Gompertz example, \texttt{rinit} sets $X_0$ to the
parameter \texttt{X0} deterministically (no process noise at
initialization), and the first observation time coincides with $t_0$, so no
process step intervenes before the first observation is assimilated: the
filtering distribution of $Z$ at the first observation time is a point mass,
$Z_1 = \log(\texttt{X0})$ with zero variance. Every subsequent observation
time is preceded by exactly one Gompertz recursion at unit time step,
matching the annual spacing of the built-in data; the \texttt{discrete\_time}
step function used by the model ignores its \texttt{dt} argument and applies
$\sigma_p$ directly as the per-step log-scale standard deviation, so no
$\sqrt{dt}$ rescaling is needed to match the recursion above exactly at
$dt=1$.

\subsection{The scalar Kalman recursion}

For $t=1,\dots,N$, with filtering mean and variance $m_{t-1},v_{t-1}$ of
$Z_{t-1}$ given $\log(\mathrm{pop}_{1:t-1})$ (and $m_1=\log(\texttt{X0})$,
$v_1=0$ at $t=1$, reflecting the deterministic initialization above):

\begin{description}
\item[Predict] (skipped at $t=1$): $m \leftarrow (1-S)\log K + S\,m$,
\quad $v \leftarrow S^2 v + \sigma_p^2$.
\item[Innovate:] forecast variance $F = v+\sigma_m^2$; innovation
$e = \log(\mathrm{pop}_t) - m$; contribution to the log likelihood
$-\tfrac12\left(\log(2\pi F) + e^2/F\right)$.
\item[Update:] Kalman gain $\kappa = v/F$; $m\leftarrow m+\kappa e$;
$v \leftarrow v - \kappa^2 F$.
\end{description}

Summing the innovation contributions over $t=1,\dots,N$ gives
$\log\mathcal{L}_{\log y}$, the exact conditional log likelihood of
$\log(\mathrm{pop}_{1:N})$ under the linear-Gaussian reduction — precisely
the computation performed by \texttt{kalman\_loglik\_logy} in
\texttt{test/gompertz\_kalman.jl}.

\subsection{The Jacobian correction}

The package's \texttt{logdmeasure} for the Gompertz model evaluates the
LogNormal density of $\mathrm{pop}_t$ itself, not of $\log(\mathrm{pop}_t)$.
Since $\mathrm{pop}_t = \exp(\log(\mathrm{pop}_t))$, the two densities are
related by the change-of-variables Jacobian
$f_{\mathrm{pop}}(y) = f_{\log\mathrm{pop}}(\log y)/y$, so, summing the log
Jacobian factor $-\log y_t$ over all observations,
\[
\log\mathcal{L}(\mathrm{pop}_{1:N}) = \log\mathcal{L}(\log(\mathrm{pop}_{1:N}))
- \sum_{t=1}^N \log(\mathrm{pop}_t).
\]
The particle filter, run with the package's actual \texttt{logdmeasure} (a
density on the natural, untransformed population scale, as it must be to
match the data), estimates $\log\mathcal{L}(\mathrm{pop}_{1:N})$ — the
left-hand side — while the Kalman recursion above computes
$\log\mathcal{L}_{\log y}=\log\mathcal{L}(\log(\mathrm{pop}_{1:N}))$, the
first term on the right. The two must therefore differ by exactly the
Jacobian sum, a fixed constant (independent of $\theta$-estimation Monte
Carlo error, and indeed independent of the particle filter's behavior
altogether) equal to \texttt{jacobian}~$=-\sum_t \log y_t$ in the test code.

\subsection{Mapping the test}

\texttt{test/gompertz\_kalman.jl} implements exactly this comparison.

\begin{lstlisting}[caption={Two checks against the exact Kalman likelihood (\texttt{test/gompertz\_kalman.jl}, lines 89--117, abridged)}]
Random.seed!(20260907)
Nps = [200,1000,5000]
nreps = 30
lme = [
    logmeanexp([logLik(pfilter(P,Np=Np,params=p1)) for _ in 1:nreps],se=true)
    for Np in Nps
]
gap = [x.est - ll_logy for x in lme]
se = [x.se for x in lme]

## primary check
tol_spread = 6*maximum(se)
@test maximum(gap)-minimum(gap) < tol_spread

## secondary check
for (g,s) in zip(gap,se)
    @test abs(g-jacobian) < 6*s
end
\end{lstlisting}

For each of three particle counts $J\in\{200,1000,5000\}$, the classic
bootstrap filter (\texttt{trigger}~$=1$, \texttt{target}~$=0$, the package
defaults) is run 30 times and the 30 log-likelihood replicates combined via
\texttt{logmeanexp} (a numerically stable $\log$ of an average of
exponentials, with a jackknife standard error) into a single estimate
\texttt{x.est} with reported Monte Carlo standard error \texttt{x.se} at
that $J$. \texttt{gap} is the difference between this estimate and the
exact log-Gompertz-scale Kalman likelihood \texttt{ll\_logy}, which by the
Jacobian identity above should equal \texttt{jacobian} up to Monte Carlo
error at every $J$ tested — \emph{not} shrink toward zero as $J$ grows.

\paragraph{Primary check.} A genuine defect in the filter or in the model
definition — as opposed to a mere additive bookkeeping error such as an
incorrectly signed Jacobian — would generically make \texttt{gap} drift
\emph{systematically with $J$}: the textbook downward bias of the log of a
Monte Carlo average of a nonnegative unbiased estimator scales as
$-\Var/(2J)$ (a Jensen's-inequality effect on $\log\widehat{\mathcal L}$, not
on $\widehat{\mathcal L}$ itself) and shrinks toward zero as $J\to\infty$; a
fixed bookkeeping constant, by contrast, shifts \texttt{gap} by the same
amount at every $J$. The primary check tests exactly this distinction:
$\max(\texttt{gap})-\min(\texttt{gap})$, the spread of \texttt{gap} across
the three particle counts, must be small relative to
\texttt{tol\_spread}~$=6\times$ the largest (noisiest, smallest-$J$) standard
error. This test is run only at the classic default settings
(\texttt{trigger}~$=1$, \texttt{target}~$=0$), where
Proposition~\ref{prop:credit}(a) already guarantees that the resampling
bookkeeping question examined in Section~\ref{sec:credit} — whether the
normalizing constant $C$ is restored at resampling, or dropped — cannot
make any difference: $C$ is identically one at $\beta=0$. This check
therefore neither tests, nor could have tested, that question one way or
the other; Section~\ref{sec:credit} settles it instead with an exact
two-particle enumeration and an end-to-end filtering test built specifically
around a latent process with memory, at intermediate $\beta$
(Section~\ref{sec:numerical-credit}), since the present check's model —
smooth and well-mixing on the log scale — is close enough to the memoryless
case that it could not have detected the bias even had this check been run
at $\beta\in(0,1)$.

\paragraph{Secondary check.} \texttt{gap} should also be close, in absolute
terms, to \texttt{jacobian} itself, at every $J$: this confirms that no fixed
bookkeeping error of the kind the primary check is designed to be
insensitive to is nonetheless actually present — i.e., that the two checks
together pin down both "no systematic drift with $J$" and "the correct fixed
constant," leaving no room for an error that happens to be constant in $J$
but wrong in value.

\section{History: a normalizing constant lost, retracted, and reinstated}
\label{sec:history}

This document did not arrive at Section~\ref{sec:credit} in one pass. The
account there was derived, implemented, then \emph{withdrawn} on the
strength of a numerical check, and then reinstated when that check was
found to have had no power. The false step is recorded here because what
it does and does not show is more instructive than the result itself.

Three candidate defects were raised while mapping each routine in
\texttt{pfilter.jl} onto its mathematical specification. Two concerned the
identity of particles rather than the likelihood.
Section~\ref{sec:resample} found that an earlier translation of the
tempered resampling step had indexed the retained weight by selection
\emph{slot} rather than by selected \emph{ancestor} --- assigning particle
$k$ the quantity $(w^k)^\beta$ instead of $(w^{p_k})^\beta$, invisible at
$\beta=0$ but targeting the wrong measure at every $\beta>0$. That defect
was real, and its author found and corrected it independently while this
document was being written. Section~\ref{sec:trace} found its cousin in
the terminal ancestry draw: a uniform draw over terminal particle indices
where those particles are not, in general, equally weighted. That one
remains uncorrected upstream.

The third is the subject of Section~\ref{sec:credit}: the normalizing
constant $C=mS'/J$ that unit-mean renormalization removes at every
tempered resampling step. It was derived, the one-line change was made ---
\texttt{systematic\_resample!} altered to return $\log C$ in place of
\texttt{nothing}, and \texttt{pfilt\_step\_comps!} to add it into the
conditional log likelihood --- and then, for a day, it was taken out
again.

\subsection{The retraction, and why it was wrong}

The argument for removing $C$ ran as follows. Proper weighting is a
\emph{sufficient} condition for an unbiased sequential Monte Carlo
likelihood estimator, not a necessary one, so the failure of that
condition does not by itself establish bias; and
Proposition~\ref{prop:credit}(c) gives $\E[C\mid w^{1:J}]=1$ exactly, so
$C$ looks like a mean-one factor that a subsequent renormalization will
absorb at no cost but variance. A paired numerical comparison seemed to
confirm it: on the Gompertz--Kalman reduction of
Section~\ref{sec:kalman}, with both estimators computed from identical
particle draws so that the difference is exactly paired, across
$\beta\in\{0.25,0.5,0.75\}$ and $J\in\{8,16,64\}$ at $400{,}000$
replicates each, the two agreed to within $0.003\%$ of
$\widehat{\mathcal L}$ with no $z$-score approaching significance.

Both halves of that were mistaken, and they failed independently.

The algebraic half fails at Section~\ref{sec:telescoping-route}. It is
true that $\E[C\mid w^{1:J}]=1$. It does not follow that
$\E[C\ell]=\E[\ell]$, because $C$ is a function of the selected ancestors
and is therefore correlated with everything those ancestors go on to
generate. A conditionally mean-one factor is harmless only if it is also
\emph{independent} of everything downstream, and this one is not. In the
two-particle case of Section~\ref{sec:numerical-credit}, $\E[C]=1$
exactly while $\operatorname{Cov}(C,\ell)=+0.025$, and
$0.875+0.025=0.9$ accounts for the entire discrepancy.

The empirical half fails for a separate reason, and it is the more
general lesson. The bias is $\Theta(1/J)$ per resampling step, and it
vanishes identically whenever the one-step predictive density is constant
across the cloud. That degeneracy is not exotic: it is exactly what
happens when the latent process has no memory, and it is nearly what
happens whenever the latent process mixes well, as the Gompertz
autoregression does. The paired design could resolve about $0.02\%$ of
$\widehat{\mathcal L}$ against a bias smaller than $0.15\%$. The null
result was a statement about the experiment, not about the estimator ---
and it was reported without the standard error that would have made that
visible. Rerunning the same comparison with the weight disparity
increased recovers the sign the theory requires.

\subsection{What each instrument can see}

The episode is a clean illustration of three distinct blind spots, and it
is worth separating them.

A comparison of \emph{average log likelihoods against an exact reference}
is blind to both defects of particle identity. Mislabelling which
ancestor's retained weight lands on which selected particle changes which
states carry which weights, and a uniform terminal draw changes which
trajectory is reported; neither changes the expected value of the sum of
weights at any observation time, which is all the marginal likelihood
ever sees. Those two were found only by reading the code against its
specification, as Sections~\ref{sec:resample} and \ref{sec:trace} do.

A comparison of \emph{one implementation against another} is blind to any
defect the two share. The grid over \texttt{trigger} and \texttt{target}
that agreed cell by cell with R \texttt{pomp}'s \texttt{wpfilter} was
read at the time as validating the tempered resampling beyond the
$\beta=0$ corner. It did not: R keeps its carried weights unnormalized
and forms each conditional log likelihood as a ratio of successive total
masses, which is algebraically the same estimator as the unit-mean
convention with $C$ discarded. Both sides dropped the same factor, so
the agreement established that the two implement the same algorithm, not
that the algorithm is unbiased.

A comparison against an \emph{exact reference on a model with no power} is
blind in the way described above. The Kalman likelihood of
Section~\ref{sec:kalman} is genuinely exact and genuinely independent of
any implementation; the model it was evaluated on simply could not
exhibit the effect at a resolvable magnitude.

What settled the question was none of these, but the exact enumeration of
Section~\ref{sec:numerical-credit}: at $J=2$ the systematic sweep is
driven by a single uniform, so the ancestry is a piecewise-constant
function of it and the expectation is a finite sum with no Monte Carlo
tolerance at all. The end-to-end test that accompanies it was then built
specifically to be \emph{capable} of exhibiting the bias --- a frozen
latent state, where the surviving particle determines the entire future,
which is the exact opposite of the memoryless degeneracy that had
concealed it.

\subsection{The lesson}

Not that term-by-term mathematical accounting is unprofitable: the same
method raised all three findings, two of which were real and the third of
which was real after all. The lessons are narrower, and both are about
how a claim is checked rather than how it is derived.

First: a null numerical result is worth nothing unless the test model is
capable of exhibiting the effect and the design has power against it. The
standard error belongs beside every null, and the question ``could this
experiment have detected the thing I am concluding is absent?'' belongs
before it.

Second: a conditionally mean-one factor is harmless only if it is also
independent of everything downstream. $\E[C]=1$ is not
$\E[C\ell]=\E[\ell]$, and the gap between them is exactly
$\operatorname{Cov}(C,\ell)$.


\end{document}
